\paragraph{De novo rare-cell discovery.}
A first line of methods detects rare populations in an unsupervised or weakly supervised setting, without requiring labeled rare examples. GiniClust \citep{jiang2016giniclust} and its successors GiniClust2/3 exploit the Gini index to identify highly variable genes expressed in small clusters. FiRE \citep{jindal2018fire} assigns a rarity score based on local density in a $k$-nearest-neighbour graph, and GapClust \citep{fang2021gapclust} identifies rare clusters by detecting gaps in the distribution of second-order density differences. CellSIUS \citep{wegmann2019cellsius} uses a sub-cluster decomposition strategy to split large consensus clusters and recover embedded rare populations. More recently, scCAD \citep{xu2024scCAD} applies anomaly detection on cluster decomposition residuals and achieves strong rare-cell identification across 25 benchmarked datasets. scSID \citep{gao2023scsid} identifies rare cell types by capturing differential expression through a sparse scoring framework. These methods excel at \emph{discovery}---identifying cell types not anticipated in advance---but they do not leverage the few labeled rare examples that are typically available in semi-supervised settings, and they do not produce annotation labels directly compatible with reference atlases.

\paragraph{Reference-based and semi-supervised annotation.}
A second line of methods assigns cell-type labels from reference atlases or training sets. SingleR \citep{aran2019reference} uses pairwise correlation against bulk reference profiles; scmap \citep{kiselev2018scmap} projects query cells onto reference centroids or $k$-nearest cells; Seurat label transfer \citep{stuart2019comprehensive} identifies shared anchors between reference and query using canonical correlation analysis; and CellTypist \citep{dominguez2022cross} trains one-versus-rest logistic regression classifiers on normalized expression profiles. scANVI \citep{xu2021probabilistic} extends scVI \citep{lopez2018deep} to the semi-supervised setting, jointly learning a batch-corrected variational autoencoder with a classifier head. scArches \citep{lotfollahi2022mapping} enables transfer learning from pre-trained reference models to new query datasets. Consensus frameworks such as popV \citep{hao2024consensus} combine multiple annotation methods to reduce per-method biases. While powerful in aggregate, these methods do not specifically address the systematic under-recovery of rare cell types under extreme label scarcity---a limitation highlighted in benchmarking studies of single-cell annotation tools \citep{abdelaal2019comparison,luecken2021benchmarking} and documented for dataset imbalance in integration workflows \citep{tian2023impacts}.

\paragraph{Imbalance-aware annotation.}
A third line explicitly targets class imbalance in single-cell annotation. sc-SynO \citep{mulvey2021scsyno} applies synthetic minority oversampling (SMOTE/LoRAS) to augment the rare-class training set before fitting a classifier. scBalance \citep{xu2023scbalance} trains a sparse neural network with adaptive weighted sampling and dropout, demonstrating improvements on 20 imbalanced scRNA-seq datasets. These methods improve rare-cell recall by modifying the training distribution, and they are complementary to scRareRefine: whereas sc-SynO and scBalance require changes to the training pipeline, scRareRefine operates post-hoc on any existing scANVI model and does not require retraining or synthetic data generation.

\paragraph{Foundation models for single-cell data.}
Large pre-trained models---including scGPT \citep{cui2024scgpt}, Geneformer \citep{theodoris2023transfer}, scFoundation \citep{hao2024large}, and CellPLM \citep{wen2023cellplm}---offer general-purpose gene and cell representations that can be fine-tuned for annotation tasks. However, zero-shot evaluation studies have shown that these models exhibit significant limitations for cell-type identification and batch robustness in the absence of task-specific fine-tuning \citep{heimberg2025zeroshot,schmidt2024evaluating}. For known rare cell types with few labeled examples, fine-tuning foundation models is itself subject to the same label imbalance problem that affects scANVI. scRareRefine is architecturally independent of the upstream embedding model and could in principle be applied downstream of any latent space that satisfies the geometric separability condition.

\paragraph{Prototype-based classification.}
Prototype networks \citep{snell2017prototypical} in the few-shot learning literature classify query examples by their distance to class centroids learned from support sets---a formulation directly analogous to our prototype scoring stage. The key differences in our setting are: (i) the embedding is pre-trained by scANVI, not jointly optimized with the prototype head; (ii) class imbalance is extreme (one or two orders of magnitude); and (iii) we use the separability ratio as an explicit gate rather than applying prototype scoring unconditionally. Closest to our approach in the single-cell context is scGAD \citep{li2022scgad}, which uses prototype representations for cross-dataset cell annotation, and SRNC \citep{li2026srnc}, which uses sequential radius-neighbor classification for imbalanced single-cell data. scRareRefine is distinguished by its explicit separability diagnostic, its multi-stage pipeline with biological verification, and its strict inductive evaluation protocol.
